% =====================================================================
% Working note
% Is the signal channel free-space, and how to add the body to h
% Companion to the coherent-exposure-operator paper (paper C)
% =====================================================================

\documentclass[11pt]{article}

\usepackage[a4paper,margin=2.6cm]{geometry}
\usepackage[utf8]{inputenc}
\usepackage[T1]{fontenc}
\usepackage{lmodern}
\usepackage{microtype}
\usepackage{amsmath,amssymb,amsthm,mathtools}
\usepackage{bm}
\usepackage{booktabs}
\usepackage{xcolor}
\usepackage[colorlinks=true,allcolors=blue!50!black]{hyperref}
\usepackage[capitalize]{cleveref}

% Macros matching paper C (build/main.tex)
\newcommand{\khat}{\hat{\bm{k}}}
\newcommand{\nhat}{\hat{\bm{n}}}
\newcommand{\rr}{\mathbf{r}}
\newcommand{\xx}{\mathbf{x}}
\newcommand{\hh}{\mathbf{h}}
\newcommand{\GG}{\mathbf{G}}
\newcommand{\Gtil}{\tilde{\mathbf{G}}}
\newcommand{\Qmat}{\mathbf{Q}}
\newcommand{\Jmat}{\mathbf{J}}
\newcommand{\Phimat}{\mathbf{\Phi}}
\newcommand{\Psimat}{\bm{\Psi}}
\newcommand{\Psitil}{\tilde{\bm{\Psi}}}
\newcommand{\Fn}{\mathbf{F}_n}
\newcommand{\Rn}{\mathbf{R}_n}
\newcommand{\psin}{\bm{\psi}_n}
\newcommand{\diff}{\mathrm{d}}
\newcommand{\Complex}{\mathbb{C}}
\newcommand{\ntilde}{\tilde{n}}
\newcommand{\Vis}{V}
\newcommand{\rue}{\rr_{\mathrm{UE}}}
\newcommand{\hb}{\hh_{\mathrm{b}}}
\DeclareMathOperator{\RE}{Re}
\DeclareMathOperator{\diag}{diag}

\theoremstyle{remark}
\newtheorem{remark}{Remark}

\title{Is the signal channel free-space, and how to add the body to $\hh$}
\author{Working note, companion to paper C}
\date{2026-08-10}

\begin{document}
\maketitle

\begin{abstract}
\noindent
The signal channel $\hh$ of paper C is the free-space channel of a
scene that does not contain the body. The body enters the model only
through the exposure branch. This note makes that statement precise,
argues why it matters for the MRT baseline and for the small
device-to-body separations of the sweep, and derives a body-perturbed
signal channel $\hb$ from the machinery the paper already has. The
perturbation has two terms, a shadowing diagonal built from the same
visibility function that gates the exposure channel, and a
physical-optics scattering term built from the Fresnel reflection
operator, the mirror image of the transmission operator $\Fn$. To
first order in the body interaction the exposure operator $\Qmat$ is
unchanged, so the consistent upgrade is: replace $\hh$ by $\hb$ in the
precoder design, keep $\Qmat$ as it is.
\end{abstract}

\section{Where the body enters the model, and where it does not}

Paper C builds three linear maps from one set of ray-traced path
triples $(j(n), \khat_n, \psin)$:
\begin{center}
\begin{tabular}{@{}llll@{}}
\toprule
Branch & Map & Amplitude block & Body present \\
\midrule
Free-space field & $\mathbf{E}(\rr) = \GG(\rr)\,\xx$,
  \ $\GG = \Psimat\,\Phimat(\rr)\,\Jmat$ & $\Psimat$ & no \\
Signal & $y = \hh^T \xx$,
  \ $\hh^T = \bm{\alpha}^T\,\Phimat(\rue)\,\Jmat$ & $\bm{\alpha}^T$ & no \\
Exposure & $S_{\mathrm{ab}}(\rr) = \|\Gtil(\rr)\,\xx\|^2$,
  \ $\Gtil = \Psitil(\rr)\,\Phimat(\rr)\,\Jmat$ & $\Psitil(\rr)$ & yes \\
\bottomrule
\end{tabular}
\end{center}
The ray trace runs in a scene without the phantom. The body is
inserted a posteriori, on the exposure branch only, as a surface that
intercepts the free-space paths and transmits them through the Fresnel
operator $\Fn(\rr)$. The signal branch samples the same free-space
field at the device position $\rue$ through the UE antenna response,
so $\hh$ is the channel of a world in which the body does not exist.
Two consequences follow.

First, the body does not shadow the device. A path arriving at $\rue$
from behind the phantom, through the torso, is counted at full
strength in $\hh$. Measurement campaigns at millimeter-wave
frequencies put human-body blockage at roughly 20 to 40\,dB, so these
are not small terms.

Second, the body does not scatter into the device. The sweep of
paper C moves the served point from 0 to 20\,cm off the chest, and
states that the device separation enters only through $\hh$. In the
model that is true because $\hh$ ignores the chest entirely. In
reality a device a few centimeters from a high-index dielectric
half-space sits in a two-ray standing wave, and its received amplitude
ripples with separation. Both effects perturb the direction of
$\hh^{*}$, which is the MRT precoder, which is the baseline against
which ECBF is scored, and which sets the exposure
$\xx^H \Qmat\, \xx$ itself.

\begin{remark}[The measured channel contains the body]
In a deployed system $\hh$ is not computed, it is estimated from
pilots that physically propagate past the user's body. Real CSI is
body-inclusive by construction, and real MRT is matched to the
body-inclusive channel. The free-space $\hh$ is therefore not the
conservative choice, it is a channel no receiver would ever measure.
Adding the body to $\hh$ moves the simulation toward the channel the
network actually sees.
\end{remark}

\section{Mechanisms and expected size}

\paragraph{Shadowing.}
In the LOS geometry of paper C the phantom stands behind the served
device, so environment paths that bounce off the far half of the hall
arrive at $\rue$ through the torso. Free-space MRT co-phases those
paths and spends transmit power on them. That power is wasted for the
link, and it is precisely the power that terminates in the trunk,
since the body is what blocks it. A body-aware $\hh$ zeroes these
paths, and MRT on it stops illuminating the body from behind. Some of
the exposure reduction that paper C attributes to ECBF may already be
delivered by matched filtering on the correct channel. That is a
hypothesis worth one experiment, not a claim.

\paragraph{Chest reflection.}
At normal incidence on skin the field reflection coefficient is
$\Gamma = (1 - \ntilde)/(1 + \ntilde)$, and with
$|\ntilde| \approx 4.83$ at 28\,GHz, $|\Gamma| \approx 0.66$. A device
at separation $d$ on the chest normal receives direct and
chest-reflected copies of the dominant path, with relative phase
$2 k_0 d$, so the received amplitude of that path is modulated by
$|1 + D\,\Gamma\, e^{-2 i k_0 d}|$ with $D \le 1$ the divergence
factor of the curved chest. On a flat chest the swing is $+4.4$ to
$-9.3$\,dB. Curvature reduces it, but at $d \le 3$\,cm the specular
patch is small and $D$ is close to one, so several dB of ripple in
$|\hh|$ across the separation sweep is the expectation, not the
exception. The current model shows none of it.

\section{The body-perturbed signal channel}

Both corrections are assembled from objects paper C already defines.
Write $c_n(\rr, \xx) = x_{j(n)}\, e^{-i k_0 \khat_n \cdot \rr}$ for
the per-path excitation, and let $\Sigma$ be the body surface with
outward normal $\nhat$, incidence cosine
$\mu_n = \nhat \cdot (-\khat_n)$, and
$\xi_n = \sqrt{\ntilde^2 - 1 + \mu_n^2}$, all as in the paper.

\subsection{Shadowing diagonal}

The exposure channel already carries the binary visibility
$\Vis(\rr, \khat_n) \in \{0,1\}$, one when an unobstructed ray reaches
$\rr$ from direction $\khat_n$. Evaluate the same function at the
device, with the phantom mesh as the occluder, and collect it into
\begin{equation}\label{eq:Vue}
  \mathbf{V}_{\mathrm{UE}}
  = \diag\bigl(\Vis(\rue, \khat_1), \ldots, \Vis(\rue, \khat_N)\bigr).
\end{equation}
The shadowed signal channel is
\begin{equation}\label{eq:h-shadow}
  \hh_{\mathrm{sh}}^T
  = \bm{\alpha}^T\, \mathbf{V}_{\mathrm{UE}}\, \Phimat(\rue)\, \Jmat .
\end{equation}
This is one ray-mesh query per path against the phantom BVH, the
identical primitive that computes the ambient-occlusion gate on the
body. The binary gate overstates the shadow edge. The refinement is
the same one the exposure branch plans for its own shadow boundaries,
replace the Heaviside by a diffraction taper, knife-edge as the first
cut and the creeping-wave (Fock) coefficient where the occluding limb
is convex on the wavelength scale, since limbs sit in the Fock regime
rather than the knife-edge regime at these frequencies.

\subsection{Fresnel reflection operator}

The transmission operator $\Fn(\rr)$ of paper C has a mirror image
that the paper never needs. For a front-facing path at
$\rr \in \Sigma$, with TE and TM bases
$\hat{e}_{s,n}, \hat{e}_{p,n}$ as in the paper, the reflected
direction and the reflected TM basis are
\begin{equation}
  \khat_n^{\mathrm{r}} = \khat_n + 2 \mu_n \nhat,
  \qquad
  \hat{e}_{p,n}^{\mathrm{r}} = \hat{e}_{s,n} \times \khat_n^{\mathrm{r}},
\end{equation}
and the Fresnel amplitude reflection coefficients for the
air-to-tissue interface are
\begin{equation}\label{eq:r-coef}
  r_{s,n} = \frac{\mu_n - \xi_n}{\mu_n + \xi_n},
  \qquad
  r_{p,n} = \frac{\ntilde^2 \mu_n - \xi_n}{\ntilde^2 \mu_n + \xi_n},
\end{equation}
the exact companions of the transmission coefficients
$t_{s,n}, t_{p,n}$ of the paper, satisfying $t_{s,n} = 1 + r_{s,n}$
and $\ntilde\, t_{p,n} = 1 + r_{p,n}$. The per-path reflection
operator is then
\begin{equation}\label{eq:Rn}
  \Rn(\rr) =
  \bigl(r_{s,n}\, \hat{e}_{s,n} \hat{e}_{s,n}^{T}
  + r_{p,n}\, \hat{e}_{p,n}^{\mathrm{r}} \hat{e}_{p,n}^{T}\bigr)\,
  \Vis(\rr, \khat_n)\, H(\mu_n),
\end{equation}
rank two, zero on back-facing and self-shadowed paths, exactly the
structure of $\Fn$ with $(t_s, t_p, \hat{e}'_p)$ replaced by
$(r_s, r_p, \hat{e}^{\mathrm{r}}_p)$. Every geometric quantity in
\cref{eq:Rn} is already computed per surface point when the exposure
kernel is assembled.

\subsection{Physical-optics scattering term}

The reflected surface field of path $n$ at $\rr' \in \Sigma$ is
$\Rn(\rr')\, \psin\, c_n(\rr', \xx)$. In the tangent-plane
(Kirchhoff) approximation each lit surface element re-radiates this
field, and the body-scattered field at the device is the surface
integral
\begin{equation}\label{eq:PO}
  \mathbf{E}_{\mathrm{sc}}(\rue)
  = \sum_n \int_{\Sigma}
  \mathbf{K}(\rue, \rr')\,
  \Rn(\rr')\, \psin\,
  e^{-i k_0 \khat_n \cdot \rr'}\, \diff A' \;
  x_{j(n)},
\end{equation}
with $\mathbf{K}$ the free-space Kirchhoff propagation kernel carrying
the spherical phase $e^{-i k_0 |\rue - \rr'|}$. The integral is
linear in $\xx$ through $x_{j(n)}$, which is the only property the
precoder design needs.

The phase along $\Sigma$ is
$\varphi_n(\rr') = -k_0 \khat_n \cdot \rr' - k_0 |\rue - \rr'|$, and
stationary phase collapses \cref{eq:PO} onto the specular points, the
surface points where $\khat_n$ reflects into the direction of the
device. With $\rr^{\star}_n$ the specular point of path $n$,
$s_n = |\rue - \rr^{\star}_n|$, and
$D_n = \sqrt{\rho_1 \rho_2 / ((\rho_1 + s_n)(\rho_2 + s_n))}$ the
divergence factor of the principal curvature radii at
$\rr^{\star}_n$, the geometrical-optics evaluation is
\begin{equation}\label{eq:GO-scatter}
  \mathbf{E}_{\mathrm{sc}, n}(\rue)
  \approx D_n\, \Rn(\rr^{\star}_n)\, \psin\,
  e^{-i k_0 (\khat_n \cdot \rr^{\star}_n + s_n)}\, x_{j(n)} .
\end{equation}
Projecting on the UE antenna response and collecting per element
through the dispatch gives the scattering correction
\begin{equation}\label{eq:dh}
  \Delta\hh^{T} = \bm{\alpha}_{\mathrm{sc}}^{T}\, \Jmat,
  \qquad
  \alpha_{\mathrm{sc}, n}
  = D_n\,
  \mathbf{C}_R(\hat{\bm{s}}_n)^{H}\, \Rn(\rr^{\star}_n)\, \psin\,
  e^{-i k_0 (\khat_n \cdot \rr^{\star}_n + s_n)},
\end{equation}
with $\hat{\bm{s}}_n$ the unit vector from $\rr^{\star}_n$ to the
device. The body-perturbed signal channel is
\begin{equation}\label{eq:hb}
  \boxed{\;
  \hb^{T}
  = \bm{\alpha}^{T}\, \mathbf{V}_{\mathrm{UE}}\, \Phimat(\rue)\, \Jmat
  \;+\; \bm{\alpha}_{\mathrm{sc}}^{T}\, \Jmat .
  \;}
\end{equation}
The first term keeps the factorization of the paper with one extra
diagonal. The second term still has the amplitude-times-dispatch
shape, but its phase is referenced at the specular point rather than
at $\rue$, so it does not factor through $\Phimat(\rue)$. That is the
only structural casualty, and it is harmless, since the precoder
design consumes $\hb$ as a single $M$-vector.

\paragraph{Two-ray sanity check.}
For the dominant path at normal incidence on the chest and the device
on the chest normal at separation $d$, the specular point sits
directly behind the device, $s_n = d$,
$\khat_n \cdot \rr^{\star}_n = \khat_n \cdot \rue + d$, and
\cref{eq:dh} reduces to
$\alpha_{\mathrm{sc}, n} = D_n\, \Gamma\, e^{-2 i k_0 d}\, \alpha_n
e^{-i k_0 \khat_n \cdot \rue}$, the textbook two-ray term with
$\Gamma$ from \cref{eq:r-coef} at $\mu_n = 1$. The separation sweep
of paper C acquires exactly the ripple of \cref{eq:PO} evaluated this
way.

\section{What happens to the exposure operator}

Nothing, to first order, and that is the point. Order the model by
powers of the body interaction. The exposure branch is already first
order, each path interacts with the surface once, through $\Fn$. The
signal branch of paper C is zeroth order, the body absent. The
corrections of \cref{eq:hb} promote the signal branch to first order,
one shadowing or one reflection event per path. A body correction to
$\Qmat$ itself would require a path to interact with the body twice,
reflect off one body part and then illuminate another, before being
absorbed. On a mostly convex body that inter-reflection term is
geometrically starved, active only in concavities such as the chin
and chest, the underarm, and between the legs, and it belongs on the
same refinement list as multilayer skin and shadow-edge diffraction.
The consistent first-order theory is therefore
\begin{equation}
  \xx^{\star}
  \propto (\lambda \Qmat + \nu \mathbf{I})^{-1} \hb^{*},
  \qquad \Qmat \text{ unchanged},
\end{equation}
and likewise MRT becomes $\xx \propto \hb^{*}$. Every downstream
quantity of paper C, the Pareto sweep, the binding-restriction
analysis, the concentration statistics, is a function of
$(\hh, \Qmat)$ and runs unmodified on $(\hb, \Qmat)$.

\section{Implementation routes}

\paragraph{Closed form, from existing machinery.}
The shadowing diagonal is $N$ ray queries against the phantom BVH per
device position, the same call the ambient-occlusion gate makes. The
reflection operator \cref{eq:Rn} reuses the per-point Fresnel bases.
The specular point of path $n$ is a mirror-point search on the mesh,
cheap because the plane-wave direction is fixed, and the curvature
radii for $D_n$ come from the mesh's principal curvatures, which the
curvature refinement of the paper already estimates. Cost is
$O(N)$ per device position, negligible against the exposure kernel.

\paragraph{Body-in-scene trace, as the oracle.}
Sionna RT accepts the phantom mesh with an ITU tissue material, and a
trace of that scene returns a body-inclusive channel directly,
including multi-bounce body-environment interactions that
\cref{eq:hb} truncates. Two cautions. First, double counting: the
exposure operator must keep being fed by the body-free trace, since
$\Fn$ models the surface interaction analytically. The clean protocol
is two traces of the same scene, body-in for $\hh$, body-out for
$\Qmat$. Second, cost: the body-in trace must be repeated per posture
and position, which is precisely the per-configuration expense the
closed-form operator exists to avoid. The oracle validates
\cref{eq:hb} on a few configurations, it does not replace it in the
sweep.

\paragraph{What to measure first.}
Three numbers decide whether this matters. One, the fraction of
$\|\hh\|^2$ carried by torso-blocked paths across the sweep, which is
the shadowing budget, computable from
$\mathbf{V}_{\mathrm{UE}}$ alone with no scattering machinery. Two,
the ripple of $\|\hb\|$ over the 0 to 3\,cm separations, which is the
two-ray budget. Three, the change in
$\xx^H \Qmat\, \xx$ at fixed received signal power when MRT switches
from $\hh^{*}$ to $\hb^{*}$, which is the exposure consequence, and
whose sign is genuinely open, since shadowing steers power away from
the trunk while the chest reflection rewards paths that bounce off
the chest into the device.

\end{document}
